基准GVA的平均ARI为0.885、错分率0.042；分离减半和作业建议参数时ARI为0.411、0.363。全部155次原拟合中，155次满足各自原停止标准；新增似然检查见补充表，内部求解器状态和参数边界另行记录。各因素的实测结果见表，不从5次重复推断普遍优势。不同目标的差异、同一目标的积分误差和参数恢复误差分别评价。Poisson场景GVA三起点平均耗时8.2秒、收敛5/5；其解析期望与剖面优化的实现不再采用旧版高维局部数值优化，但新旧硬件/容差不同，不能直接用旧耗时计算加速比。MCEM所选解的每个体观测目标投影得分最大范数范围为0.001222至0.05756；小Q增量并不直接保证观测得分为零，表中的剩余似然差也须同时查看。MCEM在二元、Poisson、负二项下的方差参数RMSE依次为0.145、0.112、0.078，混合比例RMSE依次为0.029、0.026、0.034。另有4次所选起点记录过内部求解器非成功状态，与外层停止标记分开保存。有28次所选解触及预设方差或固定效应边界，参数解释受约束影响。样本量100至1000时，GVA固定效应RMSE由0.301变为0.075；序列10至50次时ARI由0.754变为0.939。方差0.1与1.5的ARI为0.970/0.662；拟合2簇与4簇为0.593/0.822，不规则时间为0.616。基准5份数据中，每个体终目标的三起点极差最大值，QML、GVA、MCEM依次为1.19e-10、4.993e-11、0.0003116；各方法按自己的目标衡量，MCEM差异还含MC波动，不能据此证明全局最优。